% !TEX program = xelatex
% 张宇考研数学预测8套卷（2027年，数学一，卷一），原试卷第4页。
\documentclass[zihao=5,fontset=windows,oneside]{exam-zh}
\examsetup{
  page/size=a4paper,
  page/show-head=false,
  page/foot-content={第;页（共;页）},
  sealline/show=false,
  question={hang=true,show-answer=false,top-sep=5pt,bottom-sep=3pt},
  choices={columns=0,column-sep=1.6em,linesep=5pt,label-pos=top-left},
  fillin={no-answer-type=none,width=19mm},
}
\linespread{1.1}
\setCJKfamilyfont{examcoverfont}[BoldFont=FandolSong-Bold.otf]{FandolSong-Regular.otf}
\providecommand{\dd}{\,\mathrm d}
\newcommand{\WorksheetTitle}[1]{%
  \begin{center}
    {\CJKfamily{examcoverfont}\bfseries\fontsize{17}{22}\selectfont 2027年张宇考研数学预测8套卷}\\[5pt]
    {\CJKfamily{examcoverfont}\fontsize{13}{18}\selectfont 数学一 · 第1套 · #1}
  \end{center}\vspace{8pt}%
}
\newcommand{\QuestionBreak}{}
\ifdefined\QuestionCrops
  \geometry{textwidth=468bp,hcentering}
  \renewcommand{\QuestionBreak}{\clearpage}
  \AtBeginDocument{\pagestyle{empty}}
\fi

\AtBeginDocument{\hypersetup{pdftitle={2027年张宇考研数学预测8套卷 · 数学一 · 第1套 · 填空题}}}
\begin{document}

\ifdefined\QuestionCrops\else
  \WorksheetTitle{填空题}
  \section{填空题：第11～16题，每小题5分，共30分。}
  \examsetup{question/bottom-sep=11pt}
\fi

\begin{question}[index=11]
$\displaystyle\lim_{x\to 0}\left(\frac{\mathrm e^x-1}{\sin x}\right)^{\cot x}=$\fillin。
\end{question}
\QuestionBreak

\begin{question}[index=12]
设连续函数 $y=f(x)$ 在区间 $[0,2]$ 上满足 $f(x)+f(2-x)=0$，\par
则 $\displaystyle\int_0^\pi f(1+\cos x)\dd x=$\fillin。
\end{question}
\QuestionBreak

\begin{question}[index=13]
曲线 $\begin{cases}x=t-\ln(1+t^2),\\ y=\arctan t\end{cases}$ 在 $t=2$ 处的曲率为\fillin。
\end{question}
\QuestionBreak

\begin{question}[index=14]
设上半球壳 $x^2+y^2+z^2=R^2\ (z\geq 0,R>0)$ 的面密度
$\rho(x,y,z)=\sqrt{x^2+y^2+z^2}$，则其对\mbox{原点}的转动惯量为\fillin。
\end{question}
\QuestionBreak

\begin{question}[index=15]
已知3阶矩阵 $\symbfit A$ 的每行元素之和均为0，$r(\symbfit A)=2$，\par
则线性方程组 $\symbfit A\symbfit x=\symbfit 0$ 的通解为\fillin。
\end{question}
\QuestionBreak

\begin{question}[index=16]
设随机变量 $X$ 和 $Y$ 均服从 $B\!\left(1,\frac12\right)$，且 $D(X+Y)=1$，\par
则 $X$ 和 $Y$ 的相关系数 $\rho=$\fillin。
\end{question}

\end{document}
